\section{Más allá del prompt: cambiar el modelo mismo}

\subsection{Métodos de fine-tuning eficientes en parámetros}

Cuando la optimización de prompt llega a su límite, es posible seguir mejorando el desempeño modificando los pesos del modelo. Pero para un LLM con cientos de miles de millones de parámetros, el fine-tuning completo (full fine-tuning) tiene un costo muy alto. Por ello los investigadores desarrollaron una serie de \textbf{métodos eficientes en parámetros} (Parameter-Efficient Fine-Tuning, PEFT).

El docente presentó varios métodos principales:

\begin{table}[H]
\centering
\begin{tabular}{lp{8cm}}
\toprule
\textbf{Método} & \textbf{Idea central} \\
\midrule
Adapter & Inserta módulos adaptadores pequeños entre las capas del transformer y solo entrena los parámetros del adaptador \\
BitFit  & Solo ajusta los parámetros de bias del modelo (una cantidad mucho menor que el total de pesos) \\
LoRA    & Mediante una descomposición de bajo rango (low-rank decomposition) agrega una matriz auxiliar junto a la matriz de pesos original, $\Delta W = BA$, y solo entrena $B$ y $A$ \\
\bottomrule
\end{tabular}
\caption{Comparación de los métodos principales de fine-tuning eficiente en parámetros}
\end{table}

\begin{importantbox}{La ventaja arquitectónica de LoRA}
LoRA (Low-Rank Adaptation) es muy popular en la industria no solo por su eficiencia de entrenamiento, sino también por su excelente \textbf{compatibilidad con la arquitectura}:

\begin{itemize}
    \item Los pesos originales del modelo $W$ permanecen sin cambios; solo es necesario almacenar la matriz auxiliar $\Delta W$, que es pequeña
    \item En la inferencia, basta con proyectar y sumar $\Delta W$ sobre $W$: $W' = W + \Delta W$
    \item Un mismo servidor puede ejecutar un modelo base y, al mismo tiempo, cargar varios adapters LoRA distintos
    \item Distintos adapters implementan funciones diferentes (por ejemplo, ``estilo de correo profesional'' frente a ``estilo de soneto shakespeariano'') y se cambian según se necesite
\end{itemize}

En cambio, si se usa full fine-tuning, cada función requiere una copia completa del modelo, por lo que el consumo de recursos crece de forma proporcional.
\end{importantbox}

\subsection{Múltiples modelos de lenguaje ``válidos''}

El docente introdujo un concepto importante: existen múltiples modelos de lenguaje igualmente ``válidos'' (valid).

Lo explicó con un ejemplo cotidiano: dado ``The storage compartment in the back of your car is called a \_\_\_\_'', un hablante de inglés estadounidense diría ``trunk'', mientras que uno de inglés británico diría ``boot''. Ambas respuestas son correctas: representan dos modelos de lenguaje distintos e igualmente válidos.

Ejemplos similares están por todas partes:
\begin{itemize}
    \item ``The ruler of the country is called the \_\_\_\_'' $\rightarrow$ king / president / premier
    \item En el estado de Nueva Jersey, en Estados Unidos, un mismo tipo de sándwich recibe distintos nombres según la región: sub / hoagie / hero
    \item La forma en que hablas con un amigo, con tus padres o con un profesor es distinta en cada caso; esto también corresponde a distintos ``submodelos de lenguaje''
\end{itemize}

\begin{knowledgebox}{Moverse entre modelos de lenguaje válidos}
Moverse (shift) entre distintos ``modelos de lenguaje válidos'' es un tema de investigación importante. Sirve a dos propósitos a la vez:
\begin{itemize}
    \item \textbf{Conversión en producto}: hacer que el modelo adopte un tono y un estilo específicos (como el tono profesional de un chatbot de servicio al cliente)
    \item \textbf{AI Safety}: asegurar que el modelo rechace solicitudes dañinas (pasar de un modelo de lenguaje ``capaz de responder cualquier pregunta'' a uno que ``se niega a responder preguntas peligrosas'')
\end{itemize}
Sea cual sea el propósito, la técnica utilizada es la misma: actualizar los pesos para cambiar la posición del modelo en el espacio de probabilidades.
\end{knowledgebox}

El docente hizo un experimento interesante con Gemini: intentó que el modelo eligiera entre ``trunk'' y ``boot''. El resultado fue que Gemini (tras un post-training extenso) fue capaz de identificar que se trataba de una pregunta ambigua y dio ambas respuestas a la vez, sin favorecer a ninguna. Esto muestra la capacidad de ``múltiples perspectivas'' que los LLM modernos adquieren tras un ajuste cuidadoso.

\subsection{Instruction Tuning (ajuste por instrucciones)}

El proceso de transformar el comportamiento de un modelo de lenguaje, de ``repetir los datos de entrenamiento'' a ``seguir instrucciones para hacer algo útil'', se llama Instruction Tuning.

En la práctica, esto consiste en construir un conjunto de datos de instrucción-respuesta, por ejemplo:

\begin{quote}
\textbf{Goal}: Get a cool sleep on a summer day. \\
\textbf{How would you accomplish this goal?} \\
(A) Put a wet towel on the pillow \quad (B) Put a pillow in the oven \\
\textbf{Answer}: (A)
\end{quote}

\begin{importantbox}{El efecto del Instruction Tuning}
Los experimentos muestran que, después del Instruction Tuning, el desempeño del modelo mejora de forma notable en diversas tareas held-out (no vistas durante el entrenamiento), y el beneficio es aún más evidente en los modelos grandes. La esencia del Instruction Tuning es enseñar al modelo a organizar y presentar el conocimiento de sus datos de entrenamiento \textbf{de la manera que los humanos consideran útil}, en lugar de limitarse a repetir literalmente esos datos.
\end{importantbox}

\subsection{RLHF y Constitutional AI}

El proceso de \textbf{RLHF (Reinforcement Learning from Human Feedback)}:
\begin{enumerate}
    \item Se hace que el modelo genere varias respuestas para un mismo prompt
    \item Los anotadores humanos ordenan las respuestas según su preferencia
    \item Se entrena un \textbf{modelo de recompensa} (Reward Model) que simula las preferencias humanas
    \item Se usa el modelo de recompensa como señal de recompensa (reward signal) del aprendizaje por refuerzo para optimizar las salidas del modelo de lenguaje
\end{enumerate}

\textbf{Constitutional AI} (propuesto por Anthropic) va un paso más allá y se pregunta: ¿de verdad necesitamos datos de preferencia humana?

\begin{knowledgebox}{La idea central de Constitutional AI}
El método de Constitutional AI consiste en:
\begin{enumerate}
    \item Redactar un conjunto explícito de reglas, una ``constitución'' (constitution), que describe el código de conducta que el modelo debe seguir
    \item Usar un LLM para evaluar si la salida de otro LLM cumple con esas reglas
    \item Actualizar los pesos del modelo según el resultado de esa evaluación
\end{enumerate}
Este método usa al propio LLM para sustituir el papel de los anotadores humanos, lo que reduce de forma considerable el costo humano de la alineación (alignment). El docente citó un fragmento de la constitution que Anthropic publicó como ejemplo.
\end{knowledgebox}

El docente resumió el mensaje central de esta sección: sin importar qué técnica se use (ingeniería de prompt, ajuste por instrucciones, RLHF o Constitutional AI), el objetivo final es siempre el mismo: \textbf{actualizar de alguna manera los pesos del modelo de lenguaje para que su comportamiento se desplace en la dirección que deseamos}.

\subsection{Resumen de la sección}

Cuando la optimización de prompt llega a su límite, existen varias formas de modificar el modelo mismo: los métodos eficientes en parámetros como LoRA son adecuados para escenarios sensibles al costo; el Instruction Tuning enseña al modelo ``cómo responder de manera útil''; RLHF y Constitutional AI usan señales de preferencia para que las salidas del modelo se ajusten más a las expectativas humanas. Lo que todas estas técnicas tienen en común es que ajustan los pesos mediante un entrenamiento continuo de predicción de la siguiente palabra (next-word prediction), eligiendo entre múltiples ``modelos de lenguaje válidos''.

